\chapter{Single-Species Transport}
\def\shorttitle{Single-Species Transport} %

\section{MT3DMS Exercise 1, Part 1: Single-Species 2-D Transport Simulation}
The purpose of this exercise is to learn the use of MODFLOW and MT3DMS for analysis of 2-D
transport of singe-species contaminants under natural-gradient and influenced hydraulic conditions. The effect
of different transport solution techniques on plume migration are examined.

\subsection{Groundwater Flow}
A sketch of the system being modeled is shown on Figure \ref{fig:single_spec_fig1}. 
The aquifer is free and drained by a river. We consider that the flow is mainly perpendicular
to the river and thus a fixed head boundary is set on the west part of the domain
south and north sides of the model are considered to be no-flow boundary as the 
natural flow is mainly oriented from west to east.
uniform thickness of 6 m that is represented by three 2 m thick layers. 
The model measures 250 m in length and 120 m in height.
Hydraulic heads of 11.0 m is assigned to the west boundary and a fixed head of 10.0 m is assigned
to the river. A uniform horizontal hydraulic
conductivity of 10 m/d is assumed, and a well is located in the south-east of the domain.

\begin{figure}[tb]
\centering
\includegraphics[height=0.9\textwidth, angle = 270]{mt3dms_ex1_3dview.eps}
\caption[Illustration of the conceptual model ]{Illustration of the conceptual model}
\label{fig:single_spec_fig1}
\end{figure}

\subsection{Solute Transport}

The contaminant source is represented by a region on the north-west of the domain 
with a constant concentration of 10 mg/L,as shown on Figure \ref{fig:single_spec_fig1}. 
Clean water enters the upstream boundary. 
Solute is free to exit to the river of the well. 
The boundary conditions along the other model
faces are implicitly set as zero mass flux. 
The transport parameters are listed in Table \ref{single_spec_table_1} and are
assumed to be uniform in the whole model.
The objective of the exercise is to determine the maximum allowable pumping rate in the
well to have a concentration in that well below 10 \micro g/L.
This will be tested with different scenarios of grid dimension and transport solution
techniques.

\begin{table}[hb]
\caption{Flow and solute transport parameters used for MT3DMS Exercise 1.}
\begin{center}
\begin{tabular}{lc} 
\toprule
Parameter                                            &         Value \\
\midrule
Flow simulation type                                 &        steady state \\
Simulation time (days)                               &  800           \\
Model \  extent\ column\ direction\ $ (m)$           &  250            \\
Model \  extent\ row\ direction\ $ (m)  $            &  120            \\
Model \  thickness\ $(m)$                            &   15            \\
Model \ top\ elevation\ $(m)$                        &   15            \\
Horizontal hydraulic conductivity\ $(m/day)$         &   10     \\ 
Porosity, $\theta$                                   &           0.3  \\
Piezometric head upstream boundary\ $(m) $           &  11        \\ 
Piezometric head downstream boundary\ $(m) $         &  10            \\ 
Longitudinal dispersivity, $\alpha_L$                &           1.0 m \\
Horizontal transverse dispersivity, $\alpha_{TH}$    &           0.1 m \\
Vertical transverse dispersivity,   $\alpha_{TV}$    &           0.01 m \\
Effective diffusion coefficient, D$^{\star}$          &             0  \\
\bottomrule

\end{tabular}
\label{single_spec_table_1}
\end{center}
\end{table}

\subsection{Setting up the flow model}

To set up the flow model proceed as follows:

\begin{itemize}

\item In iPHT3D, create a new model and save it in a new and separate directory (folder).

\item Select Parameters $\rightarrow$ 1.Model $\rightarrow$ Model to specify the model dimension ('2D')
and type ('free')
\item Select Parameters $\rightarrow$ 1.Model $\rightarrow$ Domain to specify
the domain as described in Table \ref{single_spec_table_1}. For the first simulation select a spacing of 5 m in x direction
and 2.5 m in y direction

\item Set the Top of the aquifer in Spatial Attributes $\rightarrow$ Modflow $\rightarrow$ dis.6, 
use one\_value input 15 m in 
the window and do not forget to click on the 'OK' button. 
The bottom is already set to the default value of 0, used in this model.

\item in order to position the objects in the model, import the map. This is done in Parameters $\rightarrow$ 1.Model $\rightarrow$ Map.
select the Mt3d\_ex1.png file (only png files can be imported as background map).
The map must have the same size as the domain and will be positioned directly with its corner at the model boundaries.
The map can be viewed or hidden using the Results Panel $\rightarrow$ 1.Model $\rightarrow$ Map.

\item Select Parameters $\rightarrow$ Time $\rightarrow$ and set the total Simulation Time to 800 days. Then define the step size to 25 days

\item Set the boundary conditions for heads by selecting Spatial Attributes $\rightarrow$ Modflow $\rightarrow$ bas.3 boundary conditions corresponding to
IBOUND (Modflow). Use the zone tool in the Add Zone panel to add zones of fixed head BC. For this purpose select the line tool, go into the domain and click
the first point around coordinates (0.5, 119) and the second at (0.5,0.5). This will draw a line at the inflow of the domain.
When the second point is clicked, a dialog appears that allows to set the value of the zone, put -1 (fixed head BC). The media is always 0 in
2D domains.

\item staying in the same spatial attribute, create another line of fixed head on the river. For this purpose select the polygon tool and simply click 
along the river. To end the drawing of the zone right click on the last point. The value is again -1

\item to fix heads at the boundaries, use Spatial Attributes $\rightarrow$ Modflow $\rightarrow$ bas.5 Initial heads. provide a value of 10 in the dialog box.
this value will be used everywhere as a starting head and will be used as a BC at places where the BC where fixed (-1) previously. In order to 
set the inflow boundary, draw a line at the same place as the inflow BC and attribute a value of 11 m. When the grid is corse, it is not necessary 
to put the BC and heads zones exactly at the same place. However if a thinner grid is used, it can be necessary. For this purpose, it is possible to 
select one zone, open the zone dialog, copy the coordinates and paste them in a new zone anywhere in the model.

\item for porosity value, Select Spatial Attributes $\rightarrow$ MT3DMS $\rightarrow$ btn.11 and set the value to 0.3. 

\item for this first step, do not specify the well.

\end{itemize}
save the model changes.

\subsection{Running MODFLOW}

You have now completed to set up the flow model and you are ready to run MODFLOW.
This step will provide MT3DMS with the flow field that is used to compute 
advective-dispersive transport.

\begin{itemize}

\item Select Parameters $\rightarrow$ 2.Flow $\rightarrow$ Write

\item Select Parameters $\rightarrow$ 2.Flow $\rightarrow$ Run...

\item When the calculation has finished, a window appears, verify that 'NORMAL TERMINATION OF MODFLOW' is written.

\end{itemize}

\subsection{Visualising hydraulic heads}

\begin{itemize}

\item Visualise results by selecting Results $\rightarrow$ 2.Flow $\rightarrow$ Head. The time step can be changed
in Results $\rightarrow$ 1.Model $\rightarrow$ Time.

\end{itemize}

You should get a steady-state head distribution like shown Figure \ref{fig:single_spec_ex1_heads}.

\begin{figure}[tb]
\centering
\includegraphics[height=0.9\textwidth, angle = 270]{mt3dms_ex1_heads.eps}

\caption[Hydraulic head distribution under natural-gradient conditions.]{Hydraulic head distribution under natural-gradient conditions.}
\label{fig:single_spec_ex1_heads}
\end{figure}


\subsection{Setting up the transport model}

\begin{itemize}

\item Set the boundary conditions for concentrations by selecting 
Spatial Attributes $\rightarrow$ MT3DMS $\rightarrow$ btn.12 boundary conditions (this is the ICBUND value). 
The value must be equal to 1 (= active cell) in the domain and a zone must be created (with a polygon) around the source with a value of -1
that specifies a constant concentration.

\item Define the initial concentrations by selecting 
Spatial Attributes $\rightarrow$ MT3DMS $\rightarrow$ btn.13 concentrations. Use the same coordinates as defined previously to allocate 
the initial concentration to 10.

\item Define the advection package by selecting 
Parameters $\rightarrow$ 3.Transport $\rightarrow$ Parameters. In the dialog box choose ADV, then Major Data and in the first 
choice box select TVD as a first trial.

\item Define the dispersivity to be used by selecting, in the same dialog, the DSP module, input a value of 1. for the 
longitudinal dispersivity and 0.1 for TRPT and 0.01 for TRPV

\item Run the transport model by selecting Parameters $\rightarrow$ 3.Transport $\rightarrow$ Write and then Run

\item an MsDos windows appears and you shall see rapidly numbers streaming and then the dialog closes automatically

\end{itemize}

\subsection{Breakthrough curves and contour plots}

To visualise the breakthrough curve for the contaminant at a specific
location proceed as follows:

\begin{itemize}

\item Define the observation well Spatial Attributes $\rightarrow$ OBSERVATION $\rightarrow$ obs.1 observations
use the zone tool to draw a point at any location

\item visualize the head contours in Results $\rightarrow$ 3.Transport $\rightarrow$ Tracer. You can change the Time above 
to see another time. It is possible, after selecting one time to use the wheel mouse to stream along time.

\item Visualise the breakthrough curve by selecting Results $\rightarrow$ 5.Observation $\rightarrow$ Breakthrough, then below, 
select the name of the point you just create. A dialog appears that allows you to select value and a breakthrough curves appears.
The curve can be exported to be used in a spreadsheet.

\end{itemize}

\begin{figure}[tb]
\centering
\includegraphics[height=0.9\textwidth, angle = 270]{mt3dms_ex1_btc.eps}
\caption[Concentration breakthrough curve at the observation point 1
(column 16, row 16 and layer 1).]{Concentration breakthrough curve at observation point 1 (column 16, row 16 and layer 1).}
\label{fig:single_spec_ex1_btc}
\end{figure}

\begin{figure}[tb]
\centering
\includegraphics[height=0.9\textwidth, angle = 270]{mt3dms_ex1_2dcont.eps}
\caption[Calculated concentration distribution in layer 1 at year 3 under
natural-gradient conditions.]{Calculated concentration distribution in layer 1 at year 3 under
natural-gradient conditions.}
\label{fig:mt3dms_ex1_2dcont}
\end{figure}

\subsection{Adding a well}

At the location of the well in the map, create a well in the model by going to 
Spatial Attributes $\rightarrow$ MODFLOW $\rightarrow$ wel.1 zones for well and creating a point zone. 
The input value is the discharge (negative) or recharge (positive) set at the well, in $m^3/d$.
Then Modflow files shall be written and the model run, then writes also the MT3DMS files and run it.
The concentraitons at the well can be obtained by setting an observaiton point at the same place as the well.


\subsection{Effect of Solution Techniques}

Now let us examine the effect of transport solution techniques on the accuracy of
simulation results. Export the breakthrough curve at the observation point obtained using
the TVD option to an ASCII file. Then, change the advection solution option to
Upstream Finite-Difference and MOC, respectively, and export the calculated
concentrations to additional ASCII files. Note that dispersion and sink/source terms are
always solved using the standard finite-difference method and the GCG implicit matrix
solver.
Use Excel to graph the three breakthrough curves. 
% You should get something similar to Figure 5. 
The TVD 
%(shown as diamonds) 
and MOC 
%(shown as triangles) 
solutions are
nearly identical, as both solutions have minimal numerical dispersion. The Upstream
Finite-Difference solution 
%(shown as squares)
, on the other hand, contains a significant
amount of numerical dispersion.
In addition, check the mass balance errors for each solution. TVD and FD should have
close to zero mass balance error, while MOC has a large mass balance error at the
beginning but diminishes quickly, which is typical of a particle-tracking based solution
technique.

\subsection{varying grid dimension}
With iPHT3D it is possible to vary the dimensions of the cell grid with no change in the other model parameters.
For this purpose, go into the Parameters $\rightarrow$ 1.Model $\rightarrow$ Domain and change the grid cell sizes.